\chapter{Conclusion}
This final chapter provides a conclusive synthesis of the research conducted throughout this bachelor thesis. It summarizes the development and core findings of the distributed ROS 2 computer vision pipeline implemented for the TurtleBot 4 platform and outlines promising directions for future academic and practical engineering extensions.

\section{Summary}
This thesis investigated the design and development of a modular, vision-based person-following framework within the Robot Operating System 2 (ROS 2) ecosystem using a TurtleBot 4 mobile learning platform. While the primary objective was to overcome the geometric "blindness" of traditional distance sensors by embedding semantic awareness into the robot’s perception layer, practical evaluation revealed that a reliable, continuous person-following behavior could not be successfully sustained due to severe system latencies. To achieve target detection, the modern You Only Look Once (YOLO) deep learning architecture was deployed, while synchronized stereo depth maps provided by a Luxonis OAK-D Pro camera were utilized to compute three-dimensional spatial coordinates.

A major empirical finding of this research was the significant computational constraint of embedded hardware. Initial testing demonstrated that the robot's built-in Raspberry Pi 4 was unable to process modern neural networks locally, resulting in critical frame-rate drops. To circumvent this bottleneck, a network-distributed offloading paradigm via Eclipse CycloneDDS was established, successfully shifting the heavy computational workload entirely to an external workstation with a high-performance graphics card. The complete pipeline was neatly wrapped into a highly modular and reusable ROS 2 package named \texttt{yolobot}, governed by a robust five-state Finite State Machine. 

While individual nodes for object detection and spatial calculation operated with high metric accuracy, experimental testing in a geometrically complex indoor environment highlighted the vulnerabilities of distributed architectures. High-frequency image streaming heavily saturated standard wireless networks, introducing transport latencies. This sensor lag caused the robot to overshoot its target orientation during rotation, leading to temporary target loss. Ultimately, this thesis demonstrates that while network offloading is a highly valuable, cost-effective method for rapid software prototyping and semantic data integration, real-time closed-loop tracking remains highly sensitive to network quality.

\section{Future Work}
Building upon the architecture established in this thesis, several clear pathways emerge for future development and optimization:

\begin{itemize}
	\item \textbf{Hardware-Accelerated Edge Computing}: To fully eliminate the wireless network bottlenecks and latency issues observed during testing, future iterations should focus on compiling and deploying the YOLO model directly onto the Luxonis OAK-D Pro camera's integrated Visual Processing Unit (VPU). Executing neural network inference on the edge would remove the need to transmit uncompressed image streams over Wi-Fi, dramatically reducing transport latency and stabilizing the tracking control loop.
	\item \textbf{Concurrent Exploration SLAM}: While the current navigation stack relies on a predefined, static occupancy grid map, the system could be enhanced by running simultaneous localization and mapping (SLAM) concurrently with the person-following state machine. This would allow a human operator to simply walk ahead of the TurtleBot 4 to autonomously map out entirely unknown environments on the fly.
	\item \textbf{Multi-Target Filtering and Re-Identification}: The software logic within the \texttt{yolobot} package could be expanded by incorporating tracking algorithms (such as ByteTrack or DeepSORT). This integration would enable the robot to lock onto a specific target person, preventing tracking failures or accidental target switches in populated rooms where multiple human agents cross the camera's field of view.
\end{itemize}
